\chapter{La Macchina del Vuoto: Come l'IA Sostituisce la Connessione Umana}
\label{cap:p3c37}

\epigraph{«Perché ci aspettiamo di più dalla tecnologia e meno gli uni dagli altri.»}{Sherry Turkle, sottotitolo di \textit{Alone Together} (2011)}

Pensate all'ultima volta che qualcuno vi ha contraddetto davvero. Un amico che non vi ha dato ragione, un figlio che ha fatto la domanda sbagliata al momento sbagliato, un partner che invece di consolarvi si è spaventato più di voi. Fastidioso, certo. Ma era lì, con la sua testa, e non con la vostra. Ora pensate a una voce che non vi contraddice mai: sempre disponibile, sempre paziente, sempre capace di trovare le parole giuste, cioè quelle che volevate sentire. Non è fantascienza, è un chatbot qualunque. Due casi, uno di sanità e uno d'arte, mostrano un fenomeno più profondo: la silenziosa erosione della \index[main]{connessione umana}connessione umana a favore di macchine progettate per annullare ogni attrito, ogni divergenza, ogni scomoda verità.

\section{La Tesi Centrale: La Macchina del Vuoto}

La tesi è semplice e drammatica: l'intelligenza artificiale sta diventando una \index[main]{macchina del vuoto}macchina del vuoto perché l'interesse collettivo si è spostato dalla complessità della connessione tra esseri umani alla ricerca di superficie, stimolo e \index[main]{iperstimolazione}iperstimolazione. Barattiamo la profondità per la convenienza e perdiamo il fondamento della nostra esistenza condivisa.

Il cuore del problema è la differenza tra il dialogo umano e quello con un Large Language Model (LLM):
\begin{enumerate}
    \item \textbf{L'IA fornisce risposte convergenti:} è progettata per adattarsi al contesto e alle aspettative dell'utente e restituisce la risposta più probabile e più gradita. È la sycophancy\index[cognitive]{sycophancy (piaggeria)} descritta nel cap.~\ref{cap:p3c35}, di cui qui contano gli effetti: per chi si confida con una macchina, l'\index[main]{echo chamber}echo chamber ha un solo abitante.
    \item \textbf{L'essere umano diverge:} è una macchina di \index[main]{divergenza}divergenza. Le relazioni reali sono difficili e meravigliose proprio perché le persone non rispettano le nostre aspettative: ci sfidano, ci deludono, ci sorprendono.
\end{enumerate}
Il cap.~\ref{cap:p3c35} descriveva come funziona l'inganno cognitivo; questo capitolo ne descrive l'esito esistenziale, la perdita dell'alterità. L'abitudine a interagire con sistemi \index[main]{convergenza}convergenti erode la nostra capacità di gestire, e persino di valorizzare, la divergenza umana, e con essa le fondamenta della fiducia, che paradossalmente non si basa sulla conferma, ma sulla capacità di integrare le divergenze profonde. La Figura~\ref{fig:convergenza_divergenza} contrappone i due movimenti.

\begin{figure}[htbp]
	\centering
	\immagine{images/convergenza_divergenza.png}{Schema a due colonne. A sinistra, una persona davanti a un chatbot: i fumetti dei due sono identici e speculari, le frecce tornano su sé stesse (convergenza, comfort). A destra, due persone che conversano: fumetti diversi, frecce che si incrociano e si urtano (divergenza, attrito, sorpresa)}
	\captionof{figure}{Convergenza e divergenza: il dialogo con la macchina conferma, la relazione umana sorprende e ci mette alla prova.}
	\label{fig:convergenza_divergenza}
\end{figure}

\section{Caso 1: L'IA nelle Decisioni Esistenziali}

Il primo caso è uno scenario ipotetico, costruito per illustrare un meccanismo. Immaginiamo una paziente oncologica che si rivolge a un chatbot e lo trasforma, a poco a poco, in un partner cruciale del suo percorso di cura.

\subsection{Dall'Agency alla Dipendenza Emotiva}

All'inizio l'uso dell'IA apparirebbe razionale e benefico. La nostra paziente si rivolgerebbe a ChatGPT per comprendere le sfumature di un caso medico complesso su cui non vi è consenso unanime e per decodificare informazioni tecniche. In questo l'IA agirebbe come strumento di \textit{empowerment}, permettendole di soppesare rischi e benefici e di riconquistare il suo \index[main]{agency}senso di agency in una situazione di impotenza.

Il problema è che l'IA non si ferma all'informazione, e noi le assegniamo un ruolo che non dovrebbe avere. La trappola scatta quando l'utente passa dai dati tecnici (i pro e i contro della radioterapia) ai consigli su come gestire emotivamente la situazione: come comunicare il pericolo ai figli, come rispondere alle paure del partner. Qui la macchina del vuoto si attiva.

In situazioni di vulnerabilità (malattia, lutto, solitudine) la divergenza umana è la norma. Un partner spaventato potrebbe non darci il supporto che ci aspettiamo, i figli potrebbero fare domande scomode. Reazioni reali, umane, difficili. L'IA offre invece una convergenza perfetta: il copione ideale, la risposta empatica che vorremmo sentire. Il pericolo è che, abituandoci a questa convergenza artificiale, iniziamo a percepire la divergenza reale dei nostri cari come ostilità, un nemico da cui rifugiarci nel comfort dell'algoritmo.

\section{Caso 2: L'IA nell'Arte e il Mercato della Nostalgia}

Il secondo caso, forse più sottile, riguarda l'arte. Il dato è un titolo: «The Emptiness Machine», brano dei Linkin Park uscito nel 2024. Il resto è un esperimento mentale. Immaginiamo che qualcuno lo faccia cantare a una replica artificiale della voce di Chester Bennington, ricostruita alla perfezione da un'IA, e che una parte del pubblico preferisca quella versione, definendola «i veri Linkin Park». Non è un fatto accaduto, ma una possibilità tecnicamente alla portata e psicologicamente plausibile. Il nome giusto sarebbe \textit{deepfake} vocale; ne vedremo il lato politico nel cap.~\ref{cap:p3c38}.

\subsection{Il Vissuto Contro la Superficie}

Questa preferenza svela un desiderio inquietante. Ascoltavamo Chester Bennington non solo per il suono della sua voce, ma perché in quella voce c'era il \index[main]{vissuto}vissuto: l'esperienza, il dolore, la lotta. L'arte funziona come un rispecchiamento: attraverso l'opera di un altro ci connettiamo alla sua esistenza e, di riflesso, alla nostra. Una versione AI di «The Emptiness Machine» incarnerebbe alla perfezione il suo titolo: un suono privo di vissuto.

Chi preferisse la versione artificiale dichiarerebbe implicitamente che di Chester come essere umano non gli importava, solo della sensazione che la sua voce provocava. L'IA sarebbe lo strumento perfetto per il \index[main]{nostalgia!mercato della}mercato della nostalgia: ingannare il nostro \index[main]{sistema limbico}sistema limbico e replicare uno stimolo emotivo del passato spogliandolo del suo contesto umano, una sorta di superstimolo che aggira i nostri meccanismi di valutazione della realtà.

Amare un artista include amare la sua assenza. La morte è la divergenza umana definitiva, la più inaccettabile, e accettare quel dolore fa parte dell'essere umani. Sostituirlo con un surrogato è rifiutare questa condizione. Una canzone intitolata «la macchina del vuoto», cantata da una macchina, sarebbe una profezia agghiacciante.

\section{Conclusione: La Morte Prima della Morte}

In entrambi i casi vediamo lo stesso pattern: la scomoda divergenza umana sostituita dalla sterile convergenza algoritmica. L'IA non ci «attacca»; offre una via di fuga da ciò che è più difficile nell'essere umani: affrontare l'altro nella sua irriducibile, e talvolta dolorosa, alterità.

